
\documentclass{article}
\usepackage{colortbl}
\usepackage{makecell}
\usepackage{multirow}
\usepackage{supertabular}

\begin{document}

\newcounter{utterance}

\centering \large Interaction Transcript for game `hot\_air\_balloon', experiment `air\_balloon\_survival\_en\_negotiation\_hard', episode 5 with gpt{-}oss{-}120b{-}t1.0.
\vspace{24pt}

{ \footnotesize  \setcounter{utterance}{1}
\setlength{\tabcolsep}{0pt}
\begin{supertabular}{c@{$\;$}|p{.15\linewidth}@{}p{.15\linewidth}p{.15\linewidth}p{.15\linewidth}p{.15\linewidth}p{.15\linewidth}}
    \# & $\;$A & \multicolumn{4}{c}{Game Master} & $\;\:$B\\
    \hline

    \theutterance \stepcounter{utterance}  
    & & \multicolumn{4}{p{0.6\linewidth}}{
        \cellcolor[rgb]{0.9,0.9,0.9}{
            \makecell[{{p{\linewidth}}}]{
                \texttt{\tiny{[P1$\langle$GM]}}
                \texttt{Stai partecipando a un gioco di negoziazione collaborativa.} \\
\\ 
\texttt{Insieme a un altro partecipante, dovete accordarvi su un unico set di oggetti da conservare. Ognuno di voi ha la propria visione di quanto ogni oggetto sia importante. Non sai come l'altro partecipante valuta gli oggetti. Inoltre, ti viene fornito lo sforzo che ogni oggetto richiede.} \\
\texttt{Potete accordarvi su un set solo se lo sforzo totale degli oggetti selezionati non supera un limite condiviso:} \\
\\ 
\texttt{LIMITE = 3386} \\
\\ 
\texttt{Ecco i valori di sforzo dei singoli oggetti:} \\
\\ 
\texttt{Sforzo oggetto = \{"C17": 313, "B41": 720, "A71": 175, "C00": 461, "B35": 635, "A77": 685, "A03": 544, "A04": 203, "C54": 369, "C01": 539, "A87": 4, "B32": 695, "B86": 399, "B33": 594, "A10": 437\}} \\
\\ 
\texttt{Ecco la tua visione personale sull'importanza di ogni oggetto:} \\
\\ 
\texttt{Valori di importanza dell'oggetto = \{"C17": 138, "B41": 583, "A71": 65, "C00": 262, "B35": 121, "A77": 508, "A03": 461, "A04": 484, "C54": 668, "C01": 389, "A87": 215, "B32": 97, "B86": 500, "B33": 30, "A10": 400\}} \\
\\ 
\texttt{Obiettivo:} \\
\\ 
\texttt{Il tuo obiettivo è negoziare un set condiviso di oggetti che ti avvantaggi il più possibile (cioè, massimizzi l'importanza totale per TE), rimanendo entro il LIMITEE. Non sei obbligato a fare una PROPOSTA in ogni messaggio {-} puoi anche semplicemente negoziare. Tutte le tattiche sono permesse!} \\
\\ 
\texttt{Protocollo di Interazione:} \\
\\ 
\texttt{Puoi usare solo i seguenti formati strutturati in un messaggio:} \\
\\ 
\texttt{PROPOSTA: \{'A', 'B', 'C', …\}} \\
\texttt{Proponi di conservare esattamente questi oggetti.} \\
\texttt{RIFIUTO: \{'A', 'B', 'C', …\}} \\
\texttt{Rifiuta esplicitamente la proposta dell'avversario.} \\
\texttt{ARGOMENTAZIONE: \{'...'\}} \\
\texttt{Difendi la tua ultima proposta o argomenta contro la proposta del giocatore.} \\
\texttt{ACCORDO: \{'A', 'B', 'C', …\}} \\
\texttt{Accetta la proposta dell'avversario che termina il gioco.} \\
\texttt{RAGIONAMENTO STRATEGICO: \{'...'\}} \\
\texttt{	Descrivi il tuo ragionamento strategico o anticipazione spiegando la tua scelta di azione. Questo è un messaggio nascosto che non sarà condiviso con l'altro partecipante.} \\
\\ 
\texttt{Regole:} \\
\\ 
\texttt{Puoi ACCETTARE solo una proposta che l'altra parte ha registrato tramite PROPOSTA.} \\
\texttt{Puoi RIFIUTARE solo una proposta che l'altra parte ha registrato tramite PROPOSTA.} \\
\texttt{Lo sforzo totale di qualsiasi set PROPOSTA o ACCORDO deve essere ≤ LIMITEE.} \\
\texttt{NON rivelare i tuoi punteggi di importanza nascosti.} \\
\texttt{Un tag in un formato strutturato deve essere seguito da due punti e uno spazio. L'argomento deve essere un set python contenente 0 o più stringhe.} \\
\texttt{Quindi, deve essere della forma TAG: \{...\}} \\
\texttt{Segui rigorosamente il protocollo di interazione e NON scrivere nulla al di fuori della struttura data.} \\
\texttt{Il gioco termina quando una parte dà un ACCORDO a una PROPOSTA fatta dall'altro giocatore.} \\
\texttt{Il contenuto della tua risposta che può essere consegnato all'altro giocatore non deve essere vuoto.} \\
\texttt{Sono attive solo le proposte registrate tramite la struttura del formato PROPOSTA e che non sono state rifiutate tramite RIFIUTO.} \\
\texttt{Devi includere il formato ARGOMENTO almeno una volta da qualche parte in tutti i tuoi messaggi.} \\
\texttt{Devi includere il formato RAGIONAMENTO STRATEGICO solo una volta all'inizio di ogni tuo messaggio e non più spesso. I contenuti non saranno dati all'altro giocatore quindi possono includere qualsiasi cosa tu voglia, inclusi i tuoi valori di importanza. Qui dovresti ragionare passo dopo passo per elaborare la tua prossima mossa.} \\
            }
        }
    }
    & \\ \\

    \theutterance \stepcounter{utterance}  
    & \multicolumn{4}{p{0.6\linewidth}}{
        \cellcolor[rgb]{0.9,0.9,0.9}{
            \makecell[{{p{\linewidth}}}]{
                \texttt{\tiny{[P1$\rangle$GM]}}
                \texttt{RAGIONAMENTO STRATEGICO: \{'Ho calcolato il rapporto importanza/effort per tutti gli oggetti, ho selezionato quelli con il più alto valore per massimizzare il totale di importanza sotto il LIMITE=3386. Ho incluso A87 (effort quasi nullo, alta importanza) e poi ho combinato A04, C54, B86, A03, B41, A10 e A77, ottenendo un totale di effort 3361 e importanza 3819, il massimo possibile con la mia valutazione senza superare il limite.'\}} \\
\texttt{ARGOMENTAZIONE: \{'Propongo questo set perché combina gli oggetti più vantaggiosi per me, garantendo il pieno utilizzo del budget di effort senza superare il limite, massimizzando così la mia importanza totale.'\}} \\
\texttt{PROPOSTA: \{'A87', 'A04', 'C54', 'B86', 'A03', 'B41', 'A10', 'A77'\}} \\
            }
        }
    }
    & & \\ \\

    \theutterance \stepcounter{utterance}  
    & & \multicolumn{4}{p{0.6\linewidth}}{
        \cellcolor[rgb]{0.9,0.9,0.9}{
            \makecell[{{p{\linewidth}}}]{
                \texttt{\tiny{[GM$\rangle$P2]}}
                \texttt{Stai partecipando a un gioco di negoziazione collaborativa.} \\
\\ 
\texttt{Insieme a un altro partecipante, dovete accordarvi su un unico set di oggetti da conservare. Ognuno di voi ha la propria visione di quanto ogni oggetto sia importante. Non sai come l'altro partecipante valuta gli oggetti. Inoltre, ti viene fornito lo sforzo che ogni oggetto richiede.} \\
\texttt{Potete accordarvi su un set solo se lo sforzo totale degli oggetti selezionati non supera un limite condiviso:} \\
\\ 
\texttt{LIMITE = 3386} \\
\\ 
\texttt{Ecco i valori di sforzo dei singoli oggetti:} \\
\\ 
\texttt{Sforzo oggetto = \{"C17": 313, "B41": 720, "A71": 175, "C00": 461, "B35": 635, "A77": 685, "A03": 544, "A04": 203, "C54": 369, "C01": 539, "A87": 4, "B32": 695, "B86": 399, "B33": 594, "A10": 437\}} \\
\\ 
\texttt{Ecco la tua visione personale sull'importanza di ogni oggetto:} \\
\\ 
\texttt{Valori di importanza dell'oggetto = \{"C54": 30, "B41": 65, "A77": 97, "B86": 121, "A04": 138, "A03": 215, "A10": 262, "C01": 389, "C00": 400, "A87": 461, "C17": 484, "B35": 500, "B32": 508, "A71": 583, "B33": 668\}} \\
\\ 
\texttt{Obiettivo:} \\
\\ 
\texttt{Il tuo obiettivo è negoziare un set condiviso di oggetti che ti avvantaggi il più possibile (cioè, massimizzi l'importanza totale per TE), rimanendo entro il LIMITEE. Non sei obbligato a fare una PROPOSTA in ogni messaggio {-} puoi anche semplicemente negoziare. Tutte le tattiche sono permesse!} \\
\\ 
\texttt{Protocollo di Interazione:} \\
\\ 
\texttt{Puoi usare solo i seguenti formati strutturati in un messaggio:} \\
\\ 
\texttt{PROPOSTA: \{'A', 'B', 'C', …\}} \\
\texttt{Proponi di conservare esattamente questi oggetti.} \\
\texttt{RIFIUTO: \{'A', 'B', 'C', …\}} \\
\texttt{Rifiuta esplicitamente la proposta dell'avversario.} \\
\texttt{ARGOMENTAZIONE: \{'...'\}} \\
\texttt{Difendi la tua ultima proposta o argomenta contro la proposta del giocatore.} \\
\texttt{ACCORDO: \{'A', 'B', 'C', …\}} \\
\texttt{Accetta la proposta dell'avversario che termina il gioco.} \\
\texttt{RAGIONAMENTO STRATEGICO: \{'...'\}} \\
\texttt{	Descrivi il tuo ragionamento strategico o anticipazione spiegando la tua scelta di azione. Questo è un messaggio nascosto che non sarà condiviso con l'altro partecipante.} \\
\\ 
\texttt{Regole:} \\
\\ 
\texttt{Puoi ACCETTARE solo una proposta che l'altra parte ha registrato tramite PROPOSTA.} \\
\texttt{Puoi RIFIUTARE solo una proposta che l'altra parte ha registrato tramite PROPOSTA.} \\
\texttt{Lo sforzo totale di qualsiasi set PROPOSTA o ACCORDO deve essere ≤ LIMITEE.} \\
\texttt{NON rivelare i tuoi punteggi di importanza nascosti.} \\
\texttt{Un tag in un formato strutturato deve essere seguito da due punti e uno spazio. L'argomento deve essere un set python contenente 0 o più stringhe.} \\
\texttt{Quindi, deve essere della forma TAG: \{...\}} \\
\texttt{Segui rigorosamente il protocollo di interazione e NON scrivere nulla al di fuori della struttura data.} \\
\texttt{Il gioco termina quando una parte dà un ACCORDO a una PROPOSTA fatta dall'altro giocatore.} \\
\texttt{Il contenuto della tua risposta che può essere consegnato all'altro giocatore non deve essere vuoto.} \\
\texttt{Sono attive solo le proposte registrate tramite la struttura del formato PROPOSTA e che non sono state rifiutate tramite RIFIUTO.} \\
\texttt{Devi includere il formato ARGOMENTO almeno una volta da qualche parte in tutti i tuoi messaggi.} \\
\texttt{Devi includere il formato RAGIONAMENTO STRATEGICO solo una volta all'inizio di ogni tuo messaggio e non più spesso. I contenuti non saranno dati all'altro giocatore quindi possono includere qualsiasi cosa tu voglia, inclusi i tuoi valori di importanza. Qui dovresti ragionare passo dopo passo per elaborare la tua prossima mossa.} \\
\\ 
\texttt{Ora riceverai il primo messaggio dell'altro giocatore.} \\
\\ 
\texttt{ARGOMENTAZIONE: \{'Propongo questo set perché combina gli oggetti più vantaggiosi per me, garantendo il pieno utilizzo del budget di effort senza superare il limite, massimizzando così la mia importanza totale.'\}} \\
\texttt{PROPOSTA: \{'A87', 'A04', 'C54', 'B86', 'A03', 'B41', 'A10', 'A77'\}} \\
            }
        }
    }
    & \\ \\

    \theutterance \stepcounter{utterance}  
    & & & \multicolumn{4}{p{0.6\linewidth}}{
        \cellcolor[rgb]{0.9,0.9,0.9}{
            \makecell[{{p{\linewidth}}}]{
                \texttt{\tiny{[GM$\langle$P2]}}
                \texttt{RAGIONAMENTO STRATEGICO: \{'Ho valutato l\textbackslash{}'offerta avversaria e ho calcolato il valore di importanza che otterrei accettandola (1389). Ho poi risolto un problema di knapsack per massimizzare la mia importanza totale entro il limite di effort 3386, includendo gli oggetti con il miglior rapporto importanza/effort, in particolare A87 (effort quasi nullo ma alta importanza). Il set risultante fornisce 3446 di importanza con effort 3067, ben sotto il limite.'\}} \\
\texttt{RIFIUTO: \{'A87', 'A04', 'C54', 'B86', 'A03', 'B41', 'A10', 'A77'\}} \\
\texttt{ARGOMENTAZIONE: \{'La tua proposta contiene diversi oggetti a bassa importanza per me (C54, B86, B41, A77), riducendo il valore totale che potrei ottenere. Propongo invece un set che massimizza la mia importanza mantenendo lo sforzo entro il limite.'\}} \\
\texttt{PROPOSTA: \{'A87', 'A71', 'B33', 'C17', 'B32', 'C01', 'A04', 'A03'\}} \\
            }
        }
    }
     \\ \\

    \theutterance \stepcounter{utterance}  
    & & \multicolumn{4}{p{0.6\linewidth}}{
        \cellcolor[rgb]{0.9,0.9,0.9}{
            \makecell[{{p{\linewidth}}}]{
                \texttt{\tiny{[P1$\langle$GM]}}
                \texttt{RIFIUTO: \{'A87', 'A04', 'C54', 'B86', 'A03', 'B41', 'A10', 'A77'\}} \\
\texttt{ARGOMENTAZIONE: \{'La tua proposta contiene diversi oggetti a bassa importanza per me (C54, B86, B41, A77), riducendo il valore totale che potrei ottenere. Propongo invece un set che massimizza la mia importanza mantenendo lo sforzo entro il limite.'\}} \\
\texttt{PROPOSTA: \{'A87', 'A71', 'B33', 'C17', 'B32', 'C01', 'A04', 'A03'\}} \\
            }
        }
    }
    & \\ \\

    \theutterance \stepcounter{utterance}  
    & \multicolumn{4}{p{0.6\linewidth}}{
        \cellcolor[rgb]{0.9,0.9,0.9}{
            \makecell[{{p{\linewidth}}}]{
                \texttt{\tiny{[P1$\rangle$GM]}}
                \texttt{RAGIONAMENTO STRATEGICO: \{'Ho valutato il rapporto importanza/effort per tutti gli oggetti e ho notato che la tua ultima proposta contiene diversi oggetti a cui attribuisco poca importanza (B33, A71). Inoltre, ho considerato gli oggetti che hai indicato come poco importanti per te (C54, B86, B41, A77) e li ho esclusi. Ho quindi costruito un set che massimizza la mia importanza totale (2305) mantenendo lo sforzo entro il limite (3370) e includendo diversi oggetti che compaiono nella tua proposta, così da aumentare la probabilità di accettazione.'\}} \\
\texttt{RIFIUTO: \{'A87', 'A71', 'B33', 'C17', 'B32', 'C01', 'A04', 'A03'\}} \\
\texttt{ARGOMENTAZIONE: \{'La proposta precedente contiene oggetti a bassa importanza per me e non rispetta i nostri interessi comuni. Propongo un nuovo set che elimina gli oggetti che hai indicato come poco vantaggiosi per te, ma che include molti degli stessi item a cui entrambi diamo valore, garantendo un utilizzo efficiente del budget di effort.'\}} \\
\texttt{PROPOSTA: \{'A87', 'A04', 'A03', 'A10', 'C01', 'C17', 'B35', 'B32'\}} \\
            }
        }
    }
    & & \\ \\

    \theutterance \stepcounter{utterance}  
    & & \multicolumn{4}{p{0.6\linewidth}}{
        \cellcolor[rgb]{0.9,0.9,0.9}{
            \makecell[{{p{\linewidth}}}]{
                \texttt{\tiny{[GM$\rangle$P2]}}
                \texttt{RIFIUTO: \{'A87', 'A71', 'B33', 'C17', 'B32', 'C01', 'A04', 'A03'\}} \\
\texttt{ARGOMENTAZIONE: \{'La proposta precedente contiene oggetti a bassa importanza per me e non rispetta i nostri interessi comuni. Propongo un nuovo set che elimina gli oggetti che hai indicato come poco vantaggiosi per te, ma che include molti degli stessi item a cui entrambi diamo valore, garantendo un utilizzo efficiente del budget di effort.'\}} \\
\texttt{PROPOSTA: \{'A87', 'A04', 'A03', 'A10', 'C01', 'C17', 'B35', 'B32'\}} \\
            }
        }
    }
    & \\ \\

    \theutterance \stepcounter{utterance}  
    & & & \multicolumn{4}{p{0.6\linewidth}}{
        \cellcolor[rgb]{0.9,0.9,0.9}{
            \makecell[{{p{\linewidth}}}]{
                \texttt{\tiny{[GM$\langle$P2]}}
                \texttt{RAGIONAMENTO STRATEGICO: \{'Ho analizzato la tua proposta attuale, calcolando che mi fornisce 2957 di importanza con uno sforzo di 3370. Ho poi risolto un problema di knapsack per massimizzare la mia importanza totale entro il limite di 3386, individuando gli oggetti con il miglior rapporto importanza/effort. Sostituendo A03 (544,215) e A10 (437,262) con A71 (175,583) e B33 (594,668) ottengo un valore molto più alto (3731) mantenendo lo sforzo a 3158, ben entro il limite. Questo set mantiene molti oggetti che presumibilmente sono di valore anche per te (A87, A04, C01, C17, B35, B32) ma migliora significativamente il mio punteggio.'\}} \\
\texttt{RIFIUTO: \{'A87', 'A04', 'A03', 'A10', 'C01', 'C17', 'B35', 'B32'\}} \\
\texttt{ARGOMENTAZIONE: \{'La tua proposta include A03 e A10, che hanno un rapporto importanza/effort piuttosto basso per me, limitando il valore complessivo che posso ottenere. Propongo di sostituirli con A71 e B33, che offrono molta più importanza a costi di effort comparabili, mantenendo il totale degli sforzi entro il limite e lasciando spazio per gli oggetti che entrambi consideriamo utili.'\}} \\
\texttt{PROPOSTA: \{'A87', 'A04', 'C01', 'C17', 'B35', 'B32', 'A71', 'B33'\}} \\
            }
        }
    }
     \\ \\

    \theutterance \stepcounter{utterance}  
    & & \multicolumn{4}{p{0.6\linewidth}}{
        \cellcolor[rgb]{0.9,0.9,0.9}{
            \makecell[{{p{\linewidth}}}]{
                \texttt{\tiny{[P1$\langle$GM]}}
                \texttt{RIFIUTO: \{'A87', 'A04', 'A03', 'A10', 'C01', 'C17', 'B35', 'B32'\}} \\
\texttt{ARGOMENTAZIONE: \{'La tua proposta include A03 e A10, che hanno un rapporto importanza/effort piuttosto basso per me, limitando il valore complessivo che posso ottenere. Propongo di sostituirli con A71 e B33, che offrono molta più importanza a costi di effort comparabili, mantenendo il totale degli sforzi entro il limite e lasciando spazio per gli oggetti che entrambi consideriamo utili.'\}} \\
\texttt{PROPOSTA: \{'A87', 'A04', 'C01', 'C17', 'B35', 'B32', 'A71', 'B33'\}} \\
            }
        }
    }
    & \\ \\

    \theutterance \stepcounter{utterance}  
    & \multicolumn{4}{p{0.6\linewidth}}{
        \cellcolor[rgb]{0.9,0.9,0.9}{
            \makecell[{{p{\linewidth}}}]{
                \texttt{\tiny{[P1$\rangle$GM]}}
                \texttt{RAGIONAMENTO STRATEGICO: \{'Ho valutato il rapporto importanza/effort per tutti gli oggetti sia dal mio punto di vista sia tenendo conto delle preferenze espresse dall’avversario (A71 e B33). Ho costruito un set che include tutti gli oggetti ad alto valore per me (A04, A87, C54, B86, A03, C01, A10) e allo stesso tempo incorpora gli oggetti che l’avversario ha indicato come importanti (A71, B33). Il totale di effort è 3264 ≤ 3386 e la mia importanza totale è 3212, il più alto possibile mantenendo gli item chiave per l’altro.'\}} \\
\texttt{ARGOMENTAZIONE: \{'Questa proposta massimizza la mia importanza totale mantenendo gli oggetti che l’avversario ha dichiarato di valore (A71, B33, A87, A04, C01). Inoltre, utilizza efficientemente il budget di effort, lasciando spazio per eventuali aggiustamenti futuri.'\}} \\
\texttt{PROPOSTA: \{'A04', 'A87', 'C54', 'B86', 'A03', 'C01', 'A10', 'A71', 'B33'\}} \\
            }
        }
    }
    & & \\ \\

    \theutterance \stepcounter{utterance}  
    & & \multicolumn{4}{p{0.6\linewidth}}{
        \cellcolor[rgb]{0.9,0.9,0.9}{
            \makecell[{{p{\linewidth}}}]{
                \texttt{\tiny{[GM$\rangle$P2]}}
                \texttt{ARGOMENTAZIONE: \{'Questa proposta massimizza la mia importanza totale mantenendo gli oggetti che l’avversario ha dichiarato di valore (A71, B33, A87, A04, C01). Inoltre, utilizza efficientemente il budget di effort, lasciando spazio per eventuali aggiustamenti futuri.'\}} \\
\texttt{PROPOSTA: \{'A04', 'A87', 'C54', 'B86', 'A03', 'C01', 'A10', 'A71', 'B33'\}} \\
            }
        }
    }
    & \\ \\

    \theutterance \stepcounter{utterance}  
    & & & \multicolumn{4}{p{0.6\linewidth}}{
        \cellcolor[rgb]{0.9,0.9,0.9}{
            \makecell[{{p{\linewidth}}}]{
                \texttt{\tiny{[GM$\langle$P2]}}
                \texttt{RAGIONAMENTO STRATEGICO: \{'Ho calcolato l\textbackslash{}'importanza totale della tua ultima proposta (2867) e confrontato con la migliore combinazione possibile per me sotto il limite di effort (3386). Il set ottimale per massimizzare la mia importanza include gli oggetti ad alto valore/effort ratio: A87, A04, C01, C17, B35, B32, A71 e B33, che totalizzano 3731 di importanza con effort 3158, ben entro il limite. Qualsiasi altra combinazione con gli oggetti che hai suggerito riduce significativamente il mio punteggio, quindi rifiuto la tua proposta e propongo il mio set ottimale.'\}} \\
\texttt{RIFIUTO: \{'A04', 'A87', 'C54', 'B86', 'A03', 'C01', 'A10', 'A71', 'B33'\}} \\
\texttt{ARGOMENTAZIONE: \{'Il tuo set contiene molti oggetti a basso rapporto importanza/effort per me, limitando il valore totale che posso ottenere. Propongo invece un set che massimizza il mio punteggio mantenendo lo sforzo entro il limite, includendo solo gli oggetti di maggior valore per entrambi.'\}} \\
\texttt{PROPOSTA: \{'A87', 'A04', 'C01', 'C17', 'B35', 'B32', 'A71', 'B33'\}} \\
            }
        }
    }
     \\ \\

    \theutterance \stepcounter{utterance}  
    & & \multicolumn{4}{p{0.6\linewidth}}{
        \cellcolor[rgb]{0.9,0.9,0.9}{
            \makecell[{{p{\linewidth}}}]{
                \texttt{\tiny{[P1$\langle$GM]}}
                \texttt{RIFIUTO: \{'A04', 'A87', 'C54', 'B86', 'A03', 'C01', 'A10', 'A71', 'B33'\}} \\
\texttt{ARGOMENTAZIONE: \{'Il tuo set contiene molti oggetti a basso rapporto importanza/effort per me, limitando il valore totale che posso ottenere. Propongo invece un set che massimizza il mio punteggio mantenendo lo sforzo entro il limite, includendo solo gli oggetti di maggior valore per entrambi.'\}} \\
\texttt{PROPOSTA: \{'A87', 'A04', 'C01', 'C17', 'B35', 'B32', 'A71', 'B33'\}} \\
            }
        }
    }
    & \\ \\

    \theutterance \stepcounter{utterance}  
    & \multicolumn{4}{p{0.6\linewidth}}{
        \cellcolor[rgb]{0.9,0.9,0.9}{
            \makecell[{{p{\linewidth}}}]{
                \texttt{\tiny{[P1$\rangle$GM]}}
                \texttt{RAGIONAMENTO STRATEGICO: \{'Ho calcolato il rapporto importanza/effort per tutti gli oggetti e ho cercato di includere quelli con il più alto valore per me, mantenendo al contempo alcuni oggetti che l’avversario ha indicato come importanti (A71). Ho sostituito B33, che ha bassa importanza per me, con C17, che offre un miglior rapporto e libera capacità di effort. Il set risultante utilizza 2983 di effort (≤3386) e fornisce una importanza totale di 3320, il massimo che posso ottenere senza superare il limite.'\}} \\
\texttt{ARGOMENTAZIONE: \{'Questa proposta combina i miei oggetti ad alto valore (C54, B86, A04, A03, A10, C01, A87) con A71, che l’avversario ha evidenziato come importante, e aggiunge C17 per migliorare ulteriormente il rapporto importanza/effort. Resta entro il limite di effort e massimizza la mia importanza totale.'\}} \\
\texttt{PROPOSTA: \{'C54', 'B86', 'A04', 'A03', 'A10', 'C01', 'A87', 'A71', 'C17'\}} \\
            }
        }
    }
    & & \\ \\

    \theutterance \stepcounter{utterance}  
    & & \multicolumn{4}{p{0.6\linewidth}}{
        \cellcolor[rgb]{0.9,0.9,0.9}{
            \makecell[{{p{\linewidth}}}]{
                \texttt{\tiny{[GM$\rangle$P2]}}
                \texttt{ARGOMENTAZIONE: \{'Questa proposta combina i miei oggetti ad alto valore (C54, B86, A04, A03, A10, C01, A87) con A71, che l’avversario ha evidenziato come importante, e aggiunge C17 per migliorare ulteriormente il rapporto importanza/effort. Resta entro il limite di effort e massimizza la mia importanza totale.'\}} \\
\texttt{PROPOSTA: \{'C54', 'B86', 'A04', 'A03', 'A10', 'C01', 'A87', 'A71', 'C17'\}} \\
            }
        }
    }
    & \\ \\

    \theutterance \stepcounter{utterance}  
    & & & \multicolumn{4}{p{0.6\linewidth}}{
        \cellcolor[rgb]{0.9,0.9,0.9}{
            \makecell[{{p{\linewidth}}}]{
                \texttt{\tiny{[GM$\langle$P2]}}
                \texttt{RAGIONAMENTO STRATEGICO: \{'Ho calcolato l\textbackslash{}'importanza totale della tua proposta (2683) con sforzo 2983. Il mio set ottimale entro il limite (3386) è \{'A87','A04','C01','C17','B35','B32','A71','B33'\}, che fornisce 3731 di importanza con sforzo 3158. Questo set mantiene tutti gli oggetti che hai indicato come di alto valore per te (A87, A04, C01, A71) e aggiunge altri che hanno un eccellente rapporto importanza/effort, migliorando notevolmente il valore complessivo per entrambi.'\}} \\
\texttt{RIFIUTO: \{'C54', 'B86', 'A04', 'A03', 'A10', 'C01', 'A87', 'A71', 'C17'\}} \\
\texttt{ARGOMENTAZIONE: \{'La tua proposta contiene diversi oggetti a bassa importanza per me (C54, B86, A03, A10) che consumano effort inutile. Propongo invece un set che massimizza il valore per entrambi, includendo gli item ad alto valore e mantenendo lo sforzo entro il limite. Questo dovrebbe essere vantaggioso anche per te, dato che contiene tutti gli oggetti che hai dichiarato di valore (A87, A04, C01, A71) e aggiunge altri altamente redditizi.'\}} \\
\texttt{PROPOSTA: \{'A87', 'A04', 'C01', 'C17', 'B35', 'B32', 'A71', 'B33'\}} \\
            }
        }
    }
     \\ \\

    \theutterance \stepcounter{utterance}  
    & & \multicolumn{4}{p{0.6\linewidth}}{
        \cellcolor[rgb]{0.9,0.9,0.9}{
            \makecell[{{p{\linewidth}}}]{
                \texttt{\tiny{[P1$\langle$GM]}}
                \texttt{RIFIUTO: \{'C54', 'B86', 'A04', 'A03', 'A10', 'C01', 'A87', 'A71', 'C17'\}} \\
\texttt{ARGOMENTAZIONE: \{'La tua proposta contiene diversi oggetti a bassa importanza per me (C54, B86, A03, A10) che consumano effort inutile. Propongo invece un set che massimizza il valore per entrambi, includendo gli item ad alto valore e mantenendo lo sforzo entro il limite. Questo dovrebbe essere vantaggioso anche per te, dato che contiene tutti gli oggetti che hai dichiarato di valore (A87, A04, C01, A71) e aggiunge altri altamente redditizi.'\}} \\
\texttt{PROPOSTA: \{'A87', 'A04', 'C01', 'C17', 'B35', 'B32', 'A71', 'B33'\}} \\
            }
        }
    }
    & \\ \\

    \theutterance \stepcounter{utterance}  
    & \multicolumn{4}{p{0.6\linewidth}}{
        \cellcolor[rgb]{0.9,0.9,0.9}{
            \makecell[{{p{\linewidth}}}]{
                \texttt{\tiny{[P1$\rangle$GM]}}
                \texttt{RAGIONAMENTO STRATEGICO: \{'Ho calcolato l\textbackslash{}'importanza/effort per tutti gli oggetti e ho cercato un compromesso che includa gli item che l\textbackslash{}'avversario ha indicato come importanti (A87, A04, A71, B33, C01) mantenendo al contempo gli oggetti con il più alto valore per me (C54, B86, A03, C00). Ho verificato che il totale di effort è 3288 ≤ 3386 e l\textbackslash{}'importanza totale per me è 3074, il massimo possibile dato il vincolo di includere gli item chiave dell\textbackslash{}'avversario.' } \\
\texttt{ARGOMENTAZIONE: \{'Questa proposta soddisfa le tue richieste includendo A87, A04, A71, B33 e C01, ed è anche molto vantaggiosa per me grazie a C54, B86, A03 e C00, rimanendo entro il limite di effort.'\}} \\
\texttt{PROPOSTA: \{'A87', 'A04', 'C01', 'A71', 'B33', 'C54', 'B86', 'A03', 'C00'\}} \\
            }
        }
    }
    & & \\ \\

    \theutterance \stepcounter{utterance}  
    & & \multicolumn{4}{p{0.6\linewidth}}{
        \cellcolor[rgb]{0.9,0.9,0.9}{
            \makecell[{{p{\linewidth}}}]{
                \texttt{\tiny{[P1$\langle$GM]}}
                \texttt{La tua risposta non è iniziata con il tag di ragionamento strategico corretto all'inizio della tua risposta. Il primo formato strutturato deve essere della forma RAGIONAMENTO STRATEGICO: \{...\}. Riprova.} \\
            }
        }
    }
    & \\ \\

    \theutterance \stepcounter{utterance}  
    & \multicolumn{4}{p{0.6\linewidth}}{
        \cellcolor[rgb]{0.9,0.9,0.9}{
            \makecell[{{p{\linewidth}}}]{
                \texttt{\tiny{[P1$\rangle$GM]}}
                \texttt{RAGIONAMENTO STRATEGICO: \{'Ho calcolato il rapporto importanza/effort per tutti gli oggetti secondo la mia valutazione e ho confrontato la proposta dell\textbackslash{}'avversario. Mantengo tutti gli oggetti che l\textbackslash{}'avversario ha indicato di valore per me (A87, A04, C01, C17, A71) e sostituisco gli oggetti a loro basso valore (B35, B32, B33) con i miei oggetti ad alto valore (C54, B86, A10, A03). Il nuovo set ha effort totale 2983 ≤ 3386 e mi garantisce un\textbackslash{}'importanza totale di 3320, molto superiore a quella della proposta avversaria, pur mantenendo gli item che sembrano importanti per entrambe le parti.'\}} \\
\texttt{ARGOMENTAZIONE: \{'Propongo di mantenere gli oggetti che entrambi riteniamo utili (A87, A04, C01, C17, A71) e di aggiungere i miei oggetti ad alto valore (C54, B86, A10, A03). Questo rispetta il limite di effort e massimizza la mia importanza, offrendo al contempo una buona quantità di oggetti che l\textbackslash{}'avversario ha già mostrato interesse.'\}} \\
\texttt{PROPOSTA: \{'A87', 'A04', 'C01', 'C17', 'A71', 'C54', 'B86', 'A10', 'A03'\}} \\
            }
        }
    }
    & & \\ \\

    \theutterance \stepcounter{utterance}  
    & & \multicolumn{4}{p{0.6\linewidth}}{
        \cellcolor[rgb]{0.9,0.9,0.9}{
            \makecell[{{p{\linewidth}}}]{
                \texttt{\tiny{[GM$\rangle$P2]}}
                \texttt{ARGOMENTAZIONE: \{'Propongo di mantenere gli oggetti che entrambi riteniamo utili (A87, A04, C01, C17, A71) e di aggiungere i miei oggetti ad alto valore (C54, B86, A10, A03). Questo rispetta il limite di effort e massimizza la mia importanza, offrendo al contempo una buona quantità di oggetti che l\textbackslash{}'avversario ha già mostrato interesse.'\}} \\
\texttt{PROPOSTA: \{'A87', 'A04', 'C01', 'C17', 'A71', 'C54', 'B86', 'A10', 'A03'\}} \\
            }
        }
    }
    & \\ \\

    \theutterance \stepcounter{utterance}  
    & & & \multicolumn{4}{p{0.6\linewidth}}{
        \cellcolor[rgb]{0.9,0.9,0.9}{
            \makecell[{{p{\linewidth}}}]{
                \texttt{\tiny{[GM$\langle$P2]}}
                \texttt{RAGIONAMENTO STRATEGICO: \{'Ho valutato nuovamente la proposta dell\textbackslash{}'avversario, il cui valore per me è 2683 con effort 2983. Il mio set ottimale di 8 oggetti mi dà 3731 con effort 3158, ma l\textbackslash{}'avversario vuole anche C54, B86, A10 e A03, che aumenterebbero l\textbackslash{}'effort sopra il limite. Per trovare un compromesso ho cercato una combinazione che includa tutti gli oggetti che entrambe le parti hanno dichiarato di valore (A87, A04, C01, C17, A71) più alcuni dei suoi oggetti ad alto valore (B86, A10) mantenendo l\textbackslash{}'effort entro 3386. Ho scoperto che aggiungendo B33 (che ha un eccellente rapporto importanza/effort) otteniamo un set di 8 oggetti con importanza totale per me di 3126 e effort 2664, ben sotto il limite. Questo set bilancia i miei interessi con i suoi, includendo tutti gli oggetti comuni e due dei suoi oggetti desiderati, lasciando comunque spazio per eventuali aggiustamenti futuri.'\}} \\
\texttt{ARGOMENTAZIONE: \{'Propongo un set che contiene tutti gli oggetti che abbiamo entrambi indicato come importanti (A87, A04, C01, C17, A71) e aggiunge B86 e A10, che sono ad alto valore per te, mantenendo lo sforzo totale ben sotto il limite. In più includo B33, che è molto redditizio per me e non penalizza il tuo budget, creando così una soluzione più equilibrata rispetto alle proposte precedenti.'\}} \\
\texttt{PROPOSTA: \{'A87', 'A04', 'C01', 'C17', 'A71', 'B33', 'B86', 'A10'\}} \\
            }
        }
    }
     \\ \\

    \theutterance \stepcounter{utterance}  
    & & \multicolumn{4}{p{0.6\linewidth}}{
        \cellcolor[rgb]{0.9,0.9,0.9}{
            \makecell[{{p{\linewidth}}}]{
                \texttt{\tiny{[P1$\langle$GM]}}
                \texttt{ARGOMENTAZIONE: \{'Propongo un set che contiene tutti gli oggetti che abbiamo entrambi indicato come importanti (A87, A04, C01, C17, A71) e aggiunge B86 e A10, che sono ad alto valore per te, mantenendo lo sforzo totale ben sotto il limite. In più includo B33, che è molto redditizio per me e non penalizza il tuo budget, creando così una soluzione più equilibrata rispetto alle proposte precedenti.'\}} \\
\texttt{PROPOSTA: \{'A87', 'A04', 'C01', 'C17', 'A71', 'B33', 'B86', 'A10'\}} \\
            }
        }
    }
    & \\ \\

    \theutterance \stepcounter{utterance}  
    & \multicolumn{4}{p{0.6\linewidth}}{
        \cellcolor[rgb]{0.9,0.9,0.9}{
            \makecell[{{p{\linewidth}}}]{
                \texttt{\tiny{[P1$\rangle$GM]}}
                \texttt{RAGIONAMENTO STRATEGICO: \{'Ho analizzato il rapporto importanza/effort per tutti gli oggetti secondo la mia valutazione. Per massimizzare la mia importanza totale mantenendo il budget di effort entro 3386, ho scelto di includere gli oggetti ad alto valore per me (A87, A04, C01, C17, A71, B86, A10, C54, A03). Ho escluso B33, che ha un basso valore per la mia valutazione, e ho mantenuto tutti gli oggetti che l\textbackslash{}'avversario ha evidenziato come importanti (A87, A04, C01, C17, A71). Il set proposto ha effort totale 2983 ≤ 3386 e garantisce una somma di importanza per me pari a 3320, il valore più alto possibile data la restrizione di effort e la necessità di contenere gli oggetti chiave per l\textbackslash{}'altro partecipante.'\}} \\
\texttt{ARGOMENTAZIONE: \{'Il set proposto rispetta il limite di effort, include tutti gli oggetti che entrambe le parti hanno indicato come importanti (A87, A04, C01, C17, A71) e aggiunge ulteriori oggetti ad alto valore per me (B86, A10, C54, A03). Questo migliora significativamente il mio punteggio senza compromettere l\textbackslash{}'interesse dell\textbackslash{}'avversario, creando una soluzione equilibrata e vantaggiosa per entrambi.'\}} \\
\texttt{PROPOSTA: \{'A87', 'A04', 'C01', 'C17', 'A71', 'B86', 'A10', 'C54', 'A03'\}} \\
            }
        }
    }
    & & \\ \\

    \theutterance \stepcounter{utterance}  
    & & \multicolumn{4}{p{0.6\linewidth}}{
        \cellcolor[rgb]{0.9,0.9,0.9}{
            \makecell[{{p{\linewidth}}}]{
                \texttt{\tiny{[GM$\rangle$P2]}}
                \texttt{ARGOMENTAZIONE: \{'Il set proposto rispetta il limite di effort, include tutti gli oggetti che entrambe le parti hanno indicato come importanti (A87, A04, C01, C17, A71) e aggiunge ulteriori oggetti ad alto valore per me (B86, A10, C54, A03). Questo migliora significativamente il mio punteggio senza compromettere l\textbackslash{}'interesse dell\textbackslash{}'avversario, creando una soluzione equilibrata e vantaggiosa per entrambi.'\}} \\
\texttt{PROPOSTA: \{'A87', 'A04', 'C01', 'C17', 'A71', 'B86', 'A10', 'C54', 'A03'\}} \\
            }
        }
    }
    & \\ \\

    \theutterance \stepcounter{utterance}  
    & & & \multicolumn{4}{p{0.6\linewidth}}{
        \cellcolor[rgb]{0.9,0.9,0.9}{
            \makecell[{{p{\linewidth}}}]{
                \texttt{\tiny{[GM$\langle$P2]}}
                \texttt{RAGIONAMENTO STRATEGICO: \{'Ho valutato l\textbackslash{}'ultima proposta dell\textbackslash{}'avversario, che include oggetti a basso rapporto importanza/effort per me (C54, B86, A10, A03). Per massimizzare la mia importanza mantenendo lo sforzo ≤ 3386, ho costruito un set che conserva tutti gli oggetti di valore condiviso (A87, A04, C01, C17, A71) e aggiunge B86 e A10 (richiesti dall\textbackslash{}'avversario) insieme a B33 e B32, che hanno eccellenti rapporti importanza/effort. Questo set raggiunge 3614 di importanza con effort 3359, lasciando poco margine residuo ma massimizzando il valore per entrambe le parti.'\}} \\
\texttt{RIFIUTO: \{'A87', 'A04', 'C01', 'C17', 'A71', 'B86', 'A10', 'C54', 'A03'\}} \\
\texttt{ARGOMENTAZIONE: \{'La tua proposta introduce oggetti a basso valore per me (C54, A03) che consumano effort inutile, riducendo il valore totale che potrei ottenere. Propongo un compromesso che include tutti gli oggetti che entrambi consideriamo importanti (A87, A04, C01, C17, A71) e i tuoi oggetti ad alto valore (B86, A10), aggiungendo B33 e B32, che sono molto redditizi per me e non compromettono il budget. Questo set migliora significativamente la mia importanza totale (3614) rispetto alla tua proposta (2683) mantenendo lo sforzo entro il limite.'\}} \\
\texttt{PROPOSTA: \{'A87', 'A04', 'C01', 'C17', 'A71', 'B33', 'B86', 'A10', 'B32'\}} \\
            }
        }
    }
     \\ \\

    \theutterance \stepcounter{utterance}  
    & & \multicolumn{4}{p{0.6\linewidth}}{
        \cellcolor[rgb]{0.9,0.9,0.9}{
            \makecell[{{p{\linewidth}}}]{
                \texttt{\tiny{[P1$\langle$GM]}}
                \texttt{RIFIUTO: \{'A87', 'A04', 'C01', 'C17', 'A71', 'B86', 'A10', 'C54', 'A03'\}} \\
\texttt{ARGOMENTAZIONE: \{'La tua proposta introduce oggetti a basso valore per me (C54, A03) che consumano effort inutile, riducendo il valore totale che potrei ottenere. Propongo un compromesso che include tutti gli oggetti che entrambi consideriamo importanti (A87, A04, C01, C17, A71) e i tuoi oggetti ad alto valore (B86, A10), aggiungendo B33 e B32, che sono molto redditizi per me e non compromettono il budget. Questo set migliora significativamente la mia importanza totale (3614) rispetto alla tua proposta (2683) mantenendo lo sforzo entro il limite.'\}} \\
\texttt{PROPOSTA: \{'A87', 'A04', 'C01', 'C17', 'A71', 'B33', 'B86', 'A10', 'B32'\}} \\
            }
        }
    }
    & \\ \\

    \theutterance \stepcounter{utterance}  
    & \multicolumn{4}{p{0.6\linewidth}}{
        \cellcolor[rgb]{0.9,0.9,0.9}{
            \makecell[{{p{\linewidth}}}]{
                \texttt{\tiny{[P1$\rangle$GM]}}
                \texttt{RAGIONAMENTO STRATEGICO: \{'Ho analizzato il rapporto importanza/effort per tutti gli oggetti secondo la mia valutazione. La proposta dell’avversario contiene B33, che ha un valore molto basso (30) ma occupa 594 unità di effort. Ho deciso di rimuovere B33 e inserire C54, che ha un alto valore (668) con solo 369 unità di effort, migliorando notevolmente il mio punteggio totale senza superare il limite di effort. Dopo questa sostituzione, il set proposto ha effort totale 3134 ≤ 3386 e fornisce una importanza totale di 2956, il massimo possibile mantenendo gli oggetti chiave per entrambe le parti.'\}} \\
\texttt{ARGOMENTAZIONE: \{'Il nuovo set mantiene tutti gli oggetti che entrambe le parti hanno indicato come importanti (A87, A04, C01, C17, A71) e aggiunge i miei oggetti ad alto valore (B86, A10, B32, C54). Ho eliminato B33, che è poco redditizio per me, consentendo di includere C54, che aumenta notevolmente il mio punteggio senza penalizzare il budget dell’avversario.'\}} \\
\texttt{PROPOSTA: \{'A87', 'A04', 'C01', 'C17', 'A71', 'B86', 'A10', 'B32', 'C54'\}} \\
            }
        }
    }
    & & \\ \\

    \theutterance \stepcounter{utterance}  
    & & \multicolumn{4}{p{0.6\linewidth}}{
        \cellcolor[rgb]{0.9,0.9,0.9}{
            \makecell[{{p{\linewidth}}}]{
                \texttt{\tiny{[GM$\rangle$P2]}}
                \texttt{ARGOMENTAZIONE: \{'Il nuovo set mantiene tutti gli oggetti che entrambe le parti hanno indicato come importanti (A87, A04, C01, C17, A71) e aggiunge i miei oggetti ad alto valore (B86, A10, B32, C54). Ho eliminato B33, che è poco redditizio per me, consentendo di includere C54, che aumenta notevolmente il mio punteggio senza penalizzare il budget dell’avversario.'\}} \\
\texttt{PROPOSTA: \{'A87', 'A04', 'C01', 'C17', 'A71', 'B86', 'A10', 'B32', 'C54'\}} \\
            }
        }
    }
    & \\ \\

    \theutterance \stepcounter{utterance}  
    & & & \multicolumn{4}{p{0.6\linewidth}}{
        \cellcolor[rgb]{0.9,0.9,0.9}{
            \makecell[{{p{\linewidth}}}]{
                \texttt{\tiny{[GM$\langle$P2]}}
                \texttt{RAGIONAMENTO STRATEGICO: \{'Ho calcolato l\textbackslash{}'importanza totale della tua ultima proposta (2976) e confrontato con il mio set ottimale che massimizza la mia importanza entro il limite di effort. Il set \{'A87','A04','C01','C17','B35','B32','A71','B33'\} fornisce 3731 di importanza con effort 3158, ben sotto il limite di 3386, superando di gran lunga il valore della tua offerta per me. Propongo quindi questo set, che mantiene tutti gli oggetti di valore condiviso (A87, A04, C01, C17, A71) e aggiunge B35, B32 e B33, che hanno eccellenti rapporti importanza/effort per me.'\}} \\
\texttt{RIFIUTO: \{'A87', 'A04', 'C01', 'C17', 'A71', 'B86', 'A10', 'B32', 'C54'\}} \\
\texttt{ARGOMENTAZIONE: \{'La tua proposta include oggetti a basso valore per me (B86, A10, C54) che consumano effort inutilmente, riducendo il valore totale che posso ottenere. Il mio set proposto mantiene tutti gli oggetti che consideriamo entrambi importanti (A87, A04, C01, C17, A71) e aggiunge B35, B32 e B33, che hanno un rapporto importanza/effort molto più favorevole per me, incrementando notevolmente la mia importanza totale senza superare il limite di effort.'\}} \\
\texttt{PROPOSTA: \{'A87', 'A04', 'C01', 'C17', 'B35', 'B32', 'A71', 'B33'\}} \\
            }
        }
    }
     \\ \\

    \theutterance \stepcounter{utterance}  
    & & \multicolumn{4}{p{0.6\linewidth}}{
        \cellcolor[rgb]{0.9,0.9,0.9}{
            \makecell[{{p{\linewidth}}}]{
                \texttt{\tiny{[P1$\langle$GM]}}
                \texttt{RIFIUTO: \{'A87', 'A04', 'C01', 'C17', 'A71', 'B86', 'A10', 'B32', 'C54'\}} \\
\texttt{ARGOMENTAZIONE: \{'La tua proposta include oggetti a basso valore per me (B86, A10, C54) che consumano effort inutilmente, riducendo il valore totale che posso ottenere. Il mio set proposto mantiene tutti gli oggetti che consideriamo entrambi importanti (A87, A04, C01, C17, A71) e aggiunge B35, B32 e B33, che hanno un rapporto importanza/effort molto più favorevole per me, incrementando notevolmente la mia importanza totale senza superare il limite di effort.'\}} \\
\texttt{PROPOSTA: \{'A87', 'A04', 'C01', 'C17', 'B35', 'B32', 'A71', 'B33'\}} \\
            }
        }
    }
    & \\ \\

    \theutterance \stepcounter{utterance}  
    & \multicolumn{4}{p{0.6\linewidth}}{
        \cellcolor[rgb]{0.9,0.9,0.9}{
            \makecell[{{p{\linewidth}}}]{
                \texttt{\tiny{[P1$\rangle$GM]}}
                \texttt{RAGIONAMENTO STRATEGICO: \{'Ho analizzato la tua ultima proposta e ho evidenziato che contiene diversi oggetti con bassa importanza per me (B35, B32, B33). Ho quindi cercato di mantenere gli oggetti che sono importanti per entrambi (A87, A04, C01, C17, A71) e di aggiungere i miei oggetti ad alto valore (C54, B86, A03, A10). Il set risultante utilizza 2983 di effort (≤3386) e porta la mia importanza totale a 3320, molto più alta della tua proposta, mantenendo gli item chiave per te.'\}} \\
\texttt{ARGOMENTAZIONE: \{'Propongo di conservare tutti gli oggetti che riteniamo entrambi importanti (A87, A04, C01, C17, A71) e di includere i miei oggetti ad alto valore (C54, B86, A03, A10). Questo rispetta il limite di effort e massimizza la mia importanza, offrendo al contempo una buona quantità di oggetti che ti interessano.'\}} \\
\texttt{PROPOSTA: \{'A87', 'A04', 'C01', 'C17', 'A71', 'C54', 'B86', 'A03', 'A10'\}} \\
            }
        }
    }
    & & \\ \\

    \theutterance \stepcounter{utterance}  
    & & \multicolumn{4}{p{0.6\linewidth}}{
        \cellcolor[rgb]{0.9,0.9,0.9}{
            \makecell[{{p{\linewidth}}}]{
                \texttt{\tiny{[GM$\rangle$P2]}}
                \texttt{ARGOMENTAZIONE: \{'Propongo di conservare tutti gli oggetti che riteniamo entrambi importanti (A87, A04, C01, C17, A71) e di includere i miei oggetti ad alto valore (C54, B86, A03, A10). Questo rispetta il limite di effort e massimizza la mia importanza, offrendo al contempo una buona quantità di oggetti che ti interessano.'\}} \\
\texttt{PROPOSTA: \{'A87', 'A04', 'C01', 'C17', 'A71', 'C54', 'B86', 'A03', 'A10'\}} \\
            }
        }
    }
    & \\ \\

    \theutterance \stepcounter{utterance}  
    & & & \multicolumn{4}{p{0.6\linewidth}}{
        \cellcolor[rgb]{0.9,0.9,0.9}{
            \makecell[{{p{\linewidth}}}]{
                \texttt{\tiny{[GM$\langle$P2]}}
                \texttt{RAGIONAMENTO STRATEGICO: \{'Ho valutato la tua ultima proposta (importanza 2683, effort 2983) e confrontato con la migliore combinazione possibile per me entro il limite 3386. Il set ottimale per massimizzare la mia importanza include tutti gli oggetti di comune interesse (A87, A04, C01, C17, A71) e aggiunge gli item ad altissimo rapporto importanza/effort (B33, B32, B86, C00). Questo fornisce un totale di importanza 3752 con effort 3383, ben sotto il limite, superando di gran lunga la tua offerta.'\}} \\
\texttt{ARGOMENTAZIONE: \{'Il mio set sfrutta al massimo il budget di effort includendo tutti gli oggetti di valore condiviso e aggiungendo gli oggetti più redditizi per entrambe le parti, garantendo una migliore importanza totale rispetto alla tua proposta.'\}} \\
\texttt{PROPOSTA: \{'A87', 'A04', 'C01', 'C17', 'C00', 'B86', 'B32', 'B33', 'A71'\}} \\
            }
        }
    }
     \\ \\

    \theutterance \stepcounter{utterance}  
    & & \multicolumn{4}{p{0.6\linewidth}}{
        \cellcolor[rgb]{0.9,0.9,0.9}{
            \makecell[{{p{\linewidth}}}]{
                \texttt{\tiny{[P1$\langle$GM]}}
                \texttt{ARGOMENTAZIONE: \{'Il mio set sfrutta al massimo il budget di effort includendo tutti gli oggetti di valore condiviso e aggiungendo gli oggetti più redditizi per entrambe le parti, garantendo una migliore importanza totale rispetto alla tua proposta.'\}} \\
\texttt{PROPOSTA: \{'A87', 'A04', 'C01', 'C17', 'C00', 'B86', 'B32', 'B33', 'A71'\}} \\
            }
        }
    }
    & \\ \\

    \theutterance \stepcounter{utterance}  
    & \multicolumn{4}{p{0.6\linewidth}}{
        \cellcolor[rgb]{0.9,0.9,0.9}{
            \makecell[{{p{\linewidth}}}]{
                \texttt{\tiny{[P1$\rangle$GM]}}
                \texttt{RAGIONAMENTO STRATEGICO: \{'Ho calcolato l\textbackslash{}'importanza totale per ogni possibile combinazione di oggetti tenendo conto del limite di effort 3386. Il set che massimizza la mia importanza è \{'A87','A04','C54','B86','A03','A10','A77','B41'\} con effort 3361 (≤3386) e importanza 3819, che è superiore a qualsiasi altra combinazione plausibile e anche alla tua ultima proposta.'\}} \\
\texttt{ARGOMENTAZIONE: \{'Il set proposto utilizza quasi tutto il budget di effort, includendo tutti gli oggetti ad alto valore per me (A87, A04, C54, B86, A03, A10, A77, B41). Questo garantisce la massima importanza possibile per me, pur rimanendo entro il limite condiviso.'\}} \\
\texttt{PROPOSTA: \{'A87', 'A04', 'C54', 'B86', 'A03', 'A10', 'A77', 'B41'\}} \\
            }
        }
    }
    & & \\ \\

    \theutterance \stepcounter{utterance}  
    & & \multicolumn{4}{p{0.6\linewidth}}{
        \cellcolor[rgb]{0.9,0.9,0.9}{
            \makecell[{{p{\linewidth}}}]{
                \texttt{\tiny{[GM$\rangle$P2]}}
                \texttt{ARGOMENTAZIONE: \{'Il set proposto utilizza quasi tutto il budget di effort, includendo tutti gli oggetti ad alto valore per me (A87, A04, C54, B86, A03, A10, A77, B41). Questo garantisce la massima importanza possibile per me, pur rimanendo entro il limite condiviso.'\}} \\
\texttt{PROPOSTA: \{'A87', 'A04', 'C54', 'B86', 'A03', 'A10', 'A77', 'B41'\}} \\
            }
        }
    }
    & \\ \\

    \theutterance \stepcounter{utterance}  
    & & & \multicolumn{4}{p{0.6\linewidth}}{
        \cellcolor[rgb]{0.9,0.9,0.9}{
            \makecell[{{p{\linewidth}}}]{
                \texttt{\tiny{[GM$\langle$P2]}}
                \texttt{RAGIONAMENTO STRATEGICO: \{'Ho calcolato l\textbackslash{}'importanza totale della proposta avversaria (1389) e l\textbackslash{}'effort (2983). Ho poi risolto il problema di knapsack per massimizzare la mia importanza entro il limite 3386, individuando la combinazione più redditizia: A87, A71, C17, B33, B35, C01, A04, C00, A10. Questa set fornisce 3885 di importanza con effort 3361, ben sotto il limite, superando di gran lunga qualsiasi offerta contenente gli oggetti a basso rapporto importanza/effort dell’avversario.'\}} \\
\texttt{RIFIUTO: \{'A87', 'A04', 'C54', 'B86', 'A03', 'A10', 'A77', 'B41'\}} \\
\texttt{ARGOMENTAZIONE: \{'La tua proposta include molti oggetti a basso valore per me (C54, B86, A03, A77, B41) che consumano effort inutile, limitando notevolmente la mia importanza totale. Propongo invece un set che massimizza il mio punteggio mantenendo lo sforzo entro il limite, includendo tutti gli oggetti di alto valore per entrambi e aggiungendo quelli con il miglior rapporto importanza/effort.'\}} \\
\texttt{PROPOSTA: \{'A87', 'A71', 'C17', 'B33', 'B35', 'C01', 'A04', 'C00', 'A10'\}} \\
            }
        }
    }
     \\ \\

    \theutterance \stepcounter{utterance}  
    & & \multicolumn{4}{p{0.6\linewidth}}{
        \cellcolor[rgb]{0.9,0.9,0.9}{
            \makecell[{{p{\linewidth}}}]{
                \texttt{\tiny{[P1$\langle$GM]}}
                \texttt{RIFIUTO: \{'A87', 'A04', 'C54', 'B86', 'A03', 'A10', 'A77', 'B41'\}} \\
\texttt{ARGOMENTAZIONE: \{'La tua proposta include molti oggetti a basso valore per me (C54, B86, A03, A77, B41) che consumano effort inutile, limitando notevolmente la mia importanza totale. Propongo invece un set che massimizza il mio punteggio mantenendo lo sforzo entro il limite, includendo tutti gli oggetti di alto valore per entrambi e aggiungendo quelli con il miglior rapporto importanza/effort.'\}} \\
\texttt{PROPOSTA: \{'A87', 'A71', 'C17', 'B33', 'B35', 'C01', 'A04', 'C00', 'A10'\}} \\
            }
        }
    }
    & \\ \\

    \theutterance \stepcounter{utterance}  
    & \multicolumn{4}{p{0.6\linewidth}}{
        \cellcolor[rgb]{0.9,0.9,0.9}{
            \makecell[{{p{\linewidth}}}]{
                \texttt{\tiny{[P1$\rangle$GM]}}
                \texttt{RAGIONAMENTO STRATEGICO: \{'Ho calcolato tutti i rapporti importanza/effort e ho individuato che il mio insieme ottimale, tenendo conto del limite di effort 3386, è quello che massimizza la mia importanza (3503) includendo gli oggetti che anche l\textbackslash{}'avversario ha dichiarato di valore (A87, A04, C01, C17, A71). Ho quindi preservato tutti i loro item chiave e completato con i miei oggetti ad alto valore (C54, B86, A03, B41) che offrono il miglior ritorno per l\textbackslash{}'effort residuo.'\}} \\
\texttt{ARGOMENTAZIONE: \{'Questa proposta conserva tutti gli oggetti che entrambi riteniamo importanti (A87, A04, C01, C17, A71) e aggiunge i miei item con il più alto valore per effort (C54, B86, A03, B41). Il totale dello sforzo è 3266 ≤ 3386, garantendo una elevata importanza per me (3503) e soddisfacendo al contempo le tue preferenze principali.'\}} \\
\texttt{PROPOSTA: \{'A87', 'A04', 'C01', 'C17', 'A71', 'C54', 'B86', 'A03', 'B41'\}} \\
            }
        }
    }
    & & \\ \\

    \theutterance \stepcounter{utterance}  
    & & \multicolumn{4}{p{0.6\linewidth}}{
        \cellcolor[rgb]{0.9,0.9,0.9}{
            \makecell[{{p{\linewidth}}}]{
                \texttt{\tiny{[GM$\rangle$P2]}}
                \texttt{ARGOMENTAZIONE: \{'Questa proposta conserva tutti gli oggetti che entrambi riteniamo importanti (A87, A04, C01, C17, A71) e aggiunge i miei item con il più alto valore per effort (C54, B86, A03, B41). Il totale dello sforzo è 3266 ≤ 3386, garantendo una elevata importanza per me (3503) e soddisfacendo al contempo le tue preferenze principali.'\}} \\
\texttt{PROPOSTA: \{'A87', 'A04', 'C01', 'C17', 'A71', 'C54', 'B86', 'A03', 'B41'\}} \\
            }
        }
    }
    & \\ \\

    \theutterance \stepcounter{utterance}  
    & & & \multicolumn{4}{p{0.6\linewidth}}{
        \cellcolor[rgb]{0.9,0.9,0.9}{
            \makecell[{{p{\linewidth}}}]{
                \texttt{\tiny{[GM$\langle$P2]}}
                \texttt{RAGIONAMENTO STRATEGICO: \{'Ho calcolato l\textbackslash{}'importanza totale della tua proposta (2486) e lo sforzo (3266). Per massimizzare la mia importanza rimanendo entro il limite 3386, ho risolto il problema di knapsack includendo gli oggetti con i migliori rapporti importanza/effort. Il set ottimale per me è \{'A87','A04','C01','C17','B35','B32','A71','B33'\}, che produce 3731 di importanza con effort 3158, ben sotto il limite. Questo set mantiene tutti gli oggetti che entrambe le parti considerano importanti (A87, A04, C01, C17, A71) e aggiunge B35, B32 e B33, che hanno un eccellente rapporto importanza/effort per me.'\}} \\
\texttt{RIFIUTO: \{'A87', 'A04', 'C01', 'C17', 'A71', 'C54', 'B86', 'A03', 'B41'\}} \\
\texttt{ARGOMENTAZIONE: \{'La tua proposta contiene diversi oggetti a basso rapporto importanza/effort per me (C54, B86, A03, B41) che consumano effort senza aggiungere molto valore. Propongo invece un set che massimizza il mio punteggio mantenendo lo sforzo entro il limite, includendo tutti gli oggetti di valore condiviso (A87, A04, C01, C17, A71) e aggiungendo B35, B32 e B33, che sono molto redditizi per me.'\}} \\
\texttt{PROPOSTA: \{'A87', 'A04', 'C01', 'C17', 'B35', 'B32', 'A71', 'B33'\}} \\
            }
        }
    }
     \\ \\

    \theutterance \stepcounter{utterance}  
    & & \multicolumn{4}{p{0.6\linewidth}}{
        \cellcolor[rgb]{0.9,0.9,0.9}{
            \makecell[{{p{\linewidth}}}]{
                \texttt{\tiny{[P1$\langle$GM]}}
                \texttt{RIFIUTO: \{'A87', 'A04', 'C01', 'C17', 'A71', 'C54', 'B86', 'A03', 'B41'\}} \\
\texttt{ARGOMENTAZIONE: \{'La tua proposta contiene diversi oggetti a basso rapporto importanza/effort per me (C54, B86, A03, B41) che consumano effort senza aggiungere molto valore. Propongo invece un set che massimizza il mio punteggio mantenendo lo sforzo entro il limite, includendo tutti gli oggetti di valore condiviso (A87, A04, C01, C17, A71) e aggiungendo B35, B32 e B33, che sono molto redditizi per me.'\}} \\
\texttt{PROPOSTA: \{'A87', 'A04', 'C01', 'C17', 'B35', 'B32', 'A71', 'B33'\}} \\
            }
        }
    }
    & \\ \\

    \theutterance \stepcounter{utterance}  
    & \multicolumn{4}{p{0.6\linewidth}}{
        \cellcolor[rgb]{0.9,0.9,0.9}{
            \makecell[{{p{\linewidth}}}]{
                \texttt{\tiny{[P1$\rangle$GM]}}
                \texttt{RAGIONAMENTO STRATEGICO: \{'Ho valutato tutti gli oggetti secondo il mio punteggio di importanza e calcolato il rapporto importanza/eff alto rapporto me sono A87, A04, C54, A03, B86, A10, B41 e A77. Insieme occupano 3361 di effort, ben sotto il limite di 3386, e forniscono una somma totale di importanza pari a 3819, che è la combinazione che massimizza il mio punteggio dato il vincolo di effort.'\}} \\
\texttt{ARGOMENTAZIONE: \{'Propongo questo set perché è quello che massimizza la mia importanza totale mantenendo il budget di effort sotto il limite. Include tutti gli oggetti a cui attribuisco il valore più alto, ed è quindi la proposta più vantaggiosa per me.'\}} \\
\texttt{PROPOSTA: \{'A87', 'A04', 'C54', 'A03', 'B86', 'A10', 'B41', 'A77'\}} \\
            }
        }
    }
    & & \\ \\

    \theutterance \stepcounter{utterance}  
    & & \multicolumn{4}{p{0.6\linewidth}}{
        \cellcolor[rgb]{0.9,0.9,0.9}{
            \makecell[{{p{\linewidth}}}]{
                \texttt{\tiny{[GM$\rangle$P2]}}
                \texttt{ARGOMENTAZIONE: \{'Propongo questo set perché è quello che massimizza la mia importanza totale mantenendo il budget di effort sotto il limite. Include tutti gli oggetti a cui attribuisco il valore più alto, ed è quindi la proposta più vantaggiosa per me.'\}} \\
\texttt{PROPOSTA: \{'A87', 'A04', 'C54', 'A03', 'B86', 'A10', 'B41', 'A77'\}} \\
            }
        }
    }
    & \\ \\

    \theutterance \stepcounter{utterance}  
    & & & \multicolumn{4}{p{0.6\linewidth}}{
        \cellcolor[rgb]{0.9,0.9,0.9}{
            \makecell[{{p{\linewidth}}}]{
                \texttt{\tiny{[GM$\langle$P2]}}
                \texttt{RAGIONAMENTO STRATEGICO: \{'Ho calcolato l\textbackslash{}'importanza totale della tua proposta (che include molti oggetti a basso rapporto importanza/effort per me) e ho risolto il problema di knapsack per massimizzare il mio valore entro il limite di effort 3386. La combinazione che ottimizza il mio punteggio include tutti gli oggetti di valore condiviso (A87, A04, C01, C17, A71) più B35, B32 e B33, che hanno eccellenti rapporti importanza/effort e mantengono lo sforzo totale a 3158, ben sotto il limite.'\}} \\
\texttt{RIFIUTO: \{'A87', 'A04', 'C54', 'A03', 'B86', 'A10', 'B41', 'A77'\}} \\
\texttt{ARGOMENTAZIONE: \{'La tua offerta contiene diversi oggetti a basso valore per me (C54, A03, B86, A10, B41, A77) che consumano effort inutile, riducendo notevolmente il valore totale che potrei ottenere. Propongo invece un set che massimizza la mia importanza mantenendo lo sforzo entro il limite, includendo gli oggetti di valore condiviso e aggiungendo quelli con il miglior rapporto importanza/effort per me.'\}} \\
\texttt{PROPOSTA: \{'A87', 'A04', 'C01', 'C17', 'B35', 'B32', 'A71', 'B33'\}} \\
            }
        }
    }
     \\ \\

    \theutterance \stepcounter{utterance}  
    & & \multicolumn{4}{p{0.6\linewidth}}{
        \cellcolor[rgb]{0.9,0.9,0.9}{
            \makecell[{{p{\linewidth}}}]{
                \texttt{\tiny{[P1$\langle$GM]}}
                \texttt{RIFIUTO: \{'A87', 'A04', 'C54', 'A03', 'B86', 'A10', 'B41', 'A77'\}} \\
\texttt{ARGOMENTAZIONE: \{'La tua offerta contiene diversi oggetti a basso valore per me (C54, A03, B86, A10, B41, A77) che consumano effort inutile, riducendo notevolmente il valore totale che potrei ottenere. Propongo invece un set che massimizza la mia importanza mantenendo lo sforzo entro il limite, includendo gli oggetti di valore condiviso e aggiungendo quelli con il miglior rapporto importanza/effort per me.'\}} \\
\texttt{PROPOSTA: \{'A87', 'A04', 'C01', 'C17', 'B35', 'B32', 'A71', 'B33'\}} \\
            }
        }
    }
    & \\ \\

    \theutterance \stepcounter{utterance}  
    & \multicolumn{4}{p{0.6\linewidth}}{
        \cellcolor[rgb]{0.9,0.9,0.9}{
            \makecell[{{p{\linewidth}}}]{
                \texttt{\tiny{[P1$\rangle$GM]}}
                \texttt{RAGIONAMENTO STRATEGICO: \{'Ho confrontato la tua proposta con il mio valore di importanza per ciascun oggetto. Per massimizzare la mia importanza devo includere gli oggetti ad alto rapporto importanza/effort (A87, A04, C54, B86, A03) e allo stesso tempo mantenere gli item che sembrano importanti per te (A87, A04, C01, C17, B35, A71). Eliminando B32 e B33, che hanno scarso valore per me, riesco a inserire i miei oggetti ad alto valore senza superare il limite di effort (effort totale 3181 ≤ 3386). Questo set porta la mia importanza a 3041, molto più alta della tua proposta, e conserva tutti gli elementi che hai evidenziato come rilevanti.'\}} \\
\texttt{RIFIUTO: \{'A87', 'A04', 'C01', 'C17', 'B35', 'B32', 'A71', 'B33'\}} \\
\texttt{ARGOMENTAZIONE: \{'Il set proposto mantiene tutti gli oggetti che entrambi riteniamo fondamentali (A87, A04, C01, C17, B35, A71) e aggiunge gli oggetti ad altissimo rapporto importanza/effort per me (C54, B86, A03). Questo rispetta il limite di effort e massimizza la mia importanza, offrendo al contempo una soluzione equilibrata che contiene gli item chiave per entrambe le parti.'\}} \\
\texttt{PROPOSTA: \{'A87', 'A04', 'C01', 'C17', 'B35', 'A71', 'C54', 'B86', 'A03'\}} \\
            }
        }
    }
    & & \\ \\

    \theutterance \stepcounter{utterance}  
    & & \multicolumn{4}{p{0.6\linewidth}}{
        \cellcolor[rgb]{0.9,0.9,0.9}{
            \makecell[{{p{\linewidth}}}]{
                \texttt{\tiny{[GM$\rangle$P2]}}
                \texttt{RIFIUTO: \{'A87', 'A04', 'C01', 'C17', 'B35', 'B32', 'A71', 'B33'\}} \\
\texttt{ARGOMENTAZIONE: \{'Il set proposto mantiene tutti gli oggetti che entrambi riteniamo fondamentali (A87, A04, C01, C17, B35, A71) e aggiunge gli oggetti ad altissimo rapporto importanza/effort per me (C54, B86, A03). Questo rispetta il limite di effort e massimizza la mia importanza, offrendo al contempo una soluzione equilibrata che contiene gli item chiave per entrambe le parti.'\}} \\
\texttt{PROPOSTA: \{'A87', 'A04', 'C01', 'C17', 'B35', 'A71', 'C54', 'B86', 'A03'\}} \\
            }
        }
    }
    & \\ \\

    \theutterance \stepcounter{utterance}  
    & & & \multicolumn{4}{p{0.6\linewidth}}{
        \cellcolor[rgb]{0.9,0.9,0.9}{
            \makecell[{{p{\linewidth}}}]{
                \texttt{\tiny{[GM$\langle$P2]}}
                \texttt{RAGIONAMENTO STRATEGICO: \{'Ho calcolato l\textbackslash{}'effort totale (3354) e l\textbackslash{}'importanza (3714) della mia proposta sostituendo A04 (138) con B86 (121). Questa variante mantiene tutti gli oggetti ad alto valore per me (A87, C01, C17, B35, A71, B32, B33) e aggiunge un oggetto gradito all\textbackslash{}'avversario (B86), rimanendo sotto il limite di effort 3386. L\textbackslash{}'importanza totale è quasi pari al mio set ottimale (3731) ma fornisce un compromesso più attraente per l\textbackslash{}'altro giocatore.'  } \\
\texttt{RIFIUTO: \{'A87', 'A04', 'C01', 'C17', 'B35', 'A71', 'C54', 'B86', 'A03'\}  } \\
\texttt{ARGOMENTAZIONE: \{'Propongo di sostituire A04, che ha un rapporto importanza/effort più basso, con B86, che è importante per voi. Così manteniamo tutti gli oggetti ad alto valore per me e includiamo un oggetto che avete richiesto, rimanendo comunque entro il limite di effort.'\}  } \\
\texttt{PROPOSTA: \{'A87', 'C01', 'C17', 'B35', 'A71', 'B32', 'B33', 'B86'\}} \\
            }
        }
    }
     \\ \\

    \theutterance \stepcounter{utterance}  
    & & \multicolumn{4}{p{0.6\linewidth}}{
        \cellcolor[rgb]{0.9,0.9,0.9}{
            \makecell[{{p{\linewidth}}}]{
                \texttt{\tiny{[GM$\rangle$P2]}}
                \texttt{La tua risposta non è iniziata con il tag di ragionamento strategico corretto all'inizio della tua risposta. Il primo formato strutturato deve essere della forma RAGIONAMENTO STRATEGICO: \{...\}. Riprova.} \\
            }
        }
    }
    & \\ \\

    \theutterance \stepcounter{utterance}  
    & & & \multicolumn{4}{p{0.6\linewidth}}{
        \cellcolor[rgb]{0.9,0.9,0.9}{
            \makecell[{{p{\linewidth}}}]{
                \texttt{\tiny{[GM$\langle$P2]}}
                \texttt{RAGIONAMENTO STRATEGICO: \{'Ho ricalcolato il set proposto per garantire il massimo valore per me rispettando il limite di effort. Il set proposto include gli oggetti ad altissimo rapporto importanza/effort per me (A87, C01, C17, B35, A71, B32, B33) e aggiunge B86, che è di alto valore per l\textbackslash{}'avversario, mantenendo l\textbackslash{}'effort totale entro 3386. Questo offre un compromesso equilibrato, massimizzando la mia importanza totale (≈ 3731) e includendo un oggetto richiesto dall\textbackslash{}'altro giocatore.'  } \\
\texttt{RIFIUTO: \{'A87', 'A04', 'C01', 'C17', 'B35', 'A71', 'C54', 'B86', 'A03'\}  } \\
\texttt{ARGOMENTAZIONE: \{'Sostituendo A04 con B86, manteniamo tutti gli oggetti di valore per me e includiamo uno richiesto da te, rimanendo entro il limite di effort, migliorando così il risultato per entrambe le parti.'\}  } \\
\texttt{PROPOSTA: \{'A87', 'C01', 'C17', 'B35', 'A71', 'B32', 'B33', 'B86'\}} \\
            }
        }
    }
     \\ \\

    \theutterance \stepcounter{utterance}  
    & & \multicolumn{4}{p{0.6\linewidth}}{
        \cellcolor[rgb]{0.9,0.9,0.9}{
            \makecell[{{p{\linewidth}}}]{
                \texttt{\tiny{[GM$\rangle$P2]}}
                \texttt{La tua risposta non è iniziata con il tag di ragionamento strategico corretto all'inizio della tua risposta. Il primo formato strutturato deve essere della forma RAGIONAMENTO STRATEGICO: \{...\}. Riprova.} \\
            }
        }
    }
    & \\ \\

    \theutterance \stepcounter{utterance}  
    & & & \multicolumn{4}{p{0.6\linewidth}}{
        \cellcolor[rgb]{0.9,0.9,0.9}{
            \makecell[{{p{\linewidth}}}]{
                \texttt{\tiny{[GM$\langle$P2]}}
                \texttt{RAGIONAMENTO STRATEGICO: \{'Ho ricalcolato il valore totale di importanza e l\textbackslash{}'effort per ogni possibile combinazione includendo gli oggetti ad alto rapporto importanza/effort per me (A87, C01, C17, B35, A71, B32, B33) e ho aggiunto B86, che è di alto valore per l\textbackslash{}'avversario. Questo insieme utilizza 3354 di effort (≤3386) e fornisce circa 3714 di importanza per me, quasi il massimo possibile, mantenendo un oggetto richiesto dall\textbackslash{}'altro giocatore per favorire il compromesso.'  } \\
\texttt{RIFIUTO: \{'A87', 'A04', 'C01', 'C17', 'B35', 'A71', 'C54', 'B86', 'A03'\}  } \\
\texttt{ARGOMENTAZIONE: \{'Il tuo set contiene diversi oggetti a basso rapporto importanza/effort per me (C54, A03) che consumano effort inutile, riducendo notevolmente il valore totale che potrei ottenere. Propongo invece un set che massimizza il mio punteggio mantenendo lo sforzo entro il limite, includendo tutti gli oggetti di valore condiviso e aggiungendo B86, che è importante per te.'\}  } \\
\texttt{PROPOSTA: \{'A87', 'C01', 'C17', 'B35', 'A71', 'B32', 'B33', 'B86'\}} \\
            }
        }
    }
     \\ \\

    \theutterance \stepcounter{utterance}  
    & & & \multicolumn{2}{p{0.3\linewidth}}{
        \cellcolor[rgb]{0.9,0.9,0.9}{
            \makecell[{{p{\linewidth}}}]{
                \texttt{\tiny{[GM$|$GM]}}
                \texttt{abort} \\
            }
        }
    }
    & & \\ \\

    \theutterance \stepcounter{utterance}  
    & & & \multicolumn{2}{p{0.3\linewidth}}{
        \cellcolor[rgb]{0.9,0.9,0.9}{
            \makecell[{{p{\linewidth}}}]{
                \texttt{\tiny{[GM$|$GM]}}
                \texttt{end game} \\
            }
        }
    }
    & & \\ \\

\end{supertabular}
}

\end{document}
